\labsection{8}{Turn data-management results into visual evidence}{sec:lab08}

\begin{labproject}{Lab 8b visualization workflow}
\textbf{Coverage.} Chapter~\ref{chap:visualization} and the cleaned-data workflow from Lab 7.\\
\textbf{Goal.} Add two useful plots to Lab 7 and present a concise story from a built-in dataset.

\textbf{Part A --- Add visual evidence to the external-file report.}

Use the cleaned Lab 7 objects and choose two plot types that match the verified columns.

Example patterns:

\begin{lstlisting}[style=dsr]
# Example: categorical count bar chart
sex_counts <- table(titanic_cleaned$sex)
barplot(
  sex_counts,
  main = "Passenger count by sex",
  xlab = "Sex",
  ylab = "Count"
)

# Example: numeric distribution
hist(
  titanic_cleaned$fare,
  main = "Fare distribution",
  xlab = "Fare"
)
\end{lstlisting}

Base each claim on the computed counts or visible plot.

\textbf{Part B --- Storytelling with a built-in dataset.}

Any built-in dataset is acceptable. This example uses \texttt{mtcars}:

\begin{lstlisting}[style=dsr]
data(mtcars)

boxplot(
  mpg ~ cyl,
  data = mtcars,
  xlab = "Number of cylinders",
  ylab = "Miles per gallon",
  main = "Mileage by cylinder group"
)

plot(
  mtcars$wt,
  mtcars$mpg,
  xlab = "Weight",
  ylab = "Miles per gallon",
  main = "Weight versus mileage"
)
\end{lstlisting}

Structure the discussion as follows:
\begin{enumerate}
  \item Identify what each axis or group represents.
  \item Describe the strongest pattern without claiming causation.
  \item Connect the plot to a practical question a manager or analyst might ask.
  \item If a second plot supports the same observation, explain how the evidence complements the first plot.
\end{enumerate}
\end{labproject}

\begin{labdeliverable}
Submit the revised Lab 7 report with two plot types and a separate evidence-based story using a built-in dataset.
\end{labdeliverable}
